\chapter{Metodologia} \label{desenvolvimento}

Este capítulo descreve como o estudo foi conduzido. O objetivo foi comparar, sob as mesmas condições, cinco configurações aplicadas ao gradiente durante o treinamento, medindo ao mesmo tempo a proteção oferecida contra um ataque de inversão de gradientes e o custo imposto à qualidade do modelo. O adversário considerado é o descrito na Seção~\ref{sec:ameaca}: honesto, porém curioso, e não adaptativo. O procedimento completo está resumido no fluxograma da Figura~\ref{fig:fluxograma} e detalhado na Seção~\ref{sec:passos}.

\section{Base de Dados}

Foram utilizadas três bases públicas de classificação de imagens, amplamente adotadas em estudos de ataques e defesas, em ordem crescente de dificuldade: o MNIST, com dígitos manuscritos \cite{LeCun1998}; o Fashion-MNIST, com fotografias de roupas e calçados no mesmo formato \cite{Xiao2017}; e o CIFAR-10, com fotografias coloridas de objetos e animais \cite{Krizhevsky2009}. A Tabela~\ref{tab:datasets} resume suas características, o tamanho original e a quantidade de imagens efetivamente utilizada em cada execução.

\begin{table}[htb]
\centering
\caption{Bases de dados utilizadas.}
\label{tab:datasets}
\begin{tabular}{lcccccc}
\toprule
Base & Dimensão & Classes & Treino & Teste & Treino (usado) & Teste (usado) \\
\midrule
MNIST         & $28 \times 28 \times 1$ & 10 & 60.000 & 10.000 & 4.000 & 1.000 \\
Fashion-MNIST & $28 \times 28 \times 1$ & 10 & 60.000 & 10.000 & 4.000 & 1.000 \\
CIFAR-10      & $32 \times 32 \times 3$ & 10 & 50.000 & 10.000 & 4.000 & 1.000 \\
\bottomrule
\end{tabular}
\legend{Fonte: \citeonline{LeCun1998}, \citeonline{Xiao2017} e \citeonline{Krizhevsky2009}.}
\end{table}

Como o estudo exigiu milhares de treinamentos, cada execução utilizou as primeiras 4.000 imagens de treino e as primeiras 1.000 de teste, sempre as mesmas, de modo que todas as configurações fossem comparadas sobre exatamente as mesmas imagens. Esse recorte reduz a acurácia atingível e é retomado na Seção~\ref{sec:ameacas-validade}.

\section{Procedimento Experimental} \label{sec:passos}

\begin{enumerate}
  \item As imagens foram convertidas para o intervalo $[0, 1]$ e normalizadas canal a canal, subtraindo-se a média e dividindo-se pelo desvio-padrão da base, conforme a Equação~\ref{eq:normalizacao}. Os valores adotados foram $\mu = 0{,}1307$ e $\sigma = 0{,}3081$ no MNIST, $\mu = 0{,}2860$ e $\sigma = 0{,}3530$ no Fashion-MNIST e $\mu = (0{,}4914;\ 0{,}4822;\ 0{,}4465)$ e $\sigma = (0{,}2470;\ 0{,}2435;\ 0{,}2616)$ no CIFAR-10. Como a reconstrução do ataque ocorre nesse mesmo espaço normalizado, as imagens são convertidas de volta à escala de pixels antes do cálculo das métricas.

  \begin{equation} \label{eq:normalizacao}
  \tilde{x} = \frac{x - \mu}{\sigma}
  \end{equation}

  \item Três arquiteturas compactas foram implementadas, cobrindo vieses indutivos distintos:
  \begin{itemize}
    \item MLP: duas camadas totalmente conectadas de 512 e 256 neurônios com ReLU e camada de saída com 10 neurônios;
    \item CNN: duas camadas convolucionais de 32 e 64 filtros $3 \times 3$, cada uma seguida de ReLU e max-pooling $2 \times 2$, e classificador denso de 256 neurônios;
    \item ViT: patches de $4 \times 4$ pixels projetados em vetores de dimensão 128, token de classe, codificação posicional aprendida e 4 blocos Transformer com 4 cabeças de atenção.
  \end{itemize}

  \item Cinco configurações foram aplicadas ao gradiente, com parâmetros fixos em todas as combinações:
  \begin{itemize}
    \item sem defesa: gradiente inalterado, usado como referência;
    \item Top-$k$: retenção dos 10\% maiores valores absolutos, conforme a Equação~\ref{eq:topk};
    \item Random-$k$: retenção de 10\% das posições sorteadas, sem reescala, conforme a Equação~\ref{eq:randk};
    \item recorte e ruído: recorte do gradiente por norma $C = 1{,}0$ seguido de ruído gaussiano com $\sigma = 1{,}0$, inspirado no DP-SGD da Equação~\ref{eq:dpsgd}, mas sem recorte por exemplo e sem contabilização de orçamento, conforme discutido na Seção~\ref{sec:ameacas-validade};
    \item QSGD: quantização estocástica em $s = 16$ níveis, conforme a Equação~\ref{eq:qsgd}.
  \end{itemize}

  \item Em cada iteração do treinamento, o gradiente de todos os parâmetros foi concatenado em um único vetor, a defesa foi aplicada a esse vetor e somente o resultado foi usado na atualização. O protocolo simula uma rodada de um único cliente, sem agregação entre participantes, que é o cenário mais favorável ao adversário discutido na Seção~\ref{sec:fl}.

  \item Cada modelo foi treinado por uma época sobre as 4.000 imagens, com o otimizador Adam, taxa de aprendizado $\eta = 10^{-3}$ e perda de entropia cruzada. Com lotes de 32 amostras, uma época corresponde a 125 atualizações; com lotes de uma amostra, a 4.000. Ao final, registrou-se a acurácia nas 1.000 imagens de teste e salvou-se o modelo treinado.

  \item O ataque iDLG foi executado sobre um modelo representativo de cada combinação de base, arquitetura e defesa, correspondente ao treinamento com lote de 32 amostras na primeira repetição, totalizando 45 modelos. Para cada modelo, foram atacadas 30 imagens escolhidas de forma estratificada no conjunto de teste, três por classe, as mesmas para todas as defesas. O uso de imagens de teste garante que a mesma amostra seja atacada em todas as configurações, inclusive naquelas cujo treinamento sorteou outra ordem de apresentação. Em cada ataque, calculou-se o gradiente de uma única imagem, aplicou-se a mesma defesa do treinamento, inferiu-se o rótulo pelo sinal do gradiente do viés da última camada, conforme a Equação~\ref{eq:idlg-label}, e otimizou-se uma imagem sintética com L-BFGS para reproduzir o gradiente observado, conforme a Equação~\ref{eq:dlg}. Foram usadas 500 iterações no MLP e na CNN e 140 no ViT, cujo custo por iteração é muito maior.

  \item A privacidade foi medida pelo erro quadrático médio entre a imagem original e a reconstruída, conforme a Equação~\ref{eq:mse}, e pela proporção de rótulos inferidos corretamente, conforme a Equação~\ref{eq:label-acc}. Para dar escala ao erro de reconstrução, define-se um valor de referência $\mathrm{MSE}_{\text{ref}}$, calculado entre cada imagem atacada e uma imagem sorteada da mesma distribuição usada para inicializar o ataque. Esse valor indica o erro esperado quando o adversário nada recupera: defesas cujo erro se aproxima dele anulam o vazamento, enquanto valores intermediários indicam reconstrução parcial. A utilidade foi medida pela acurácia no conjunto de teste, conforme a Equação~\ref{eq:acc}.

  \item Por fim, aplicou-se o plano de análise estatística descrito na Seção~\ref{sec:analise}.
\end{enumerate}

O desenho fatorial completo, combinando três arquiteturas, três bases, cinco configurações e dois tamanhos de lote, resultou em 90 combinações. Cada combinação foi repetida 30 vezes com sementes aleatórias distintas, totalizando 2.700 treinamentos e 1.350 ataques. Todas as implementações foram feitas em Python, com as bibliotecas PyTorch e torchvision, e os modelos treinados foram salvos em disco para que o ataque fosse aplicado exatamente aos mesmos parâmetros produzidos no treinamento. O código, as sementes, os modelos treinados e os arquivos com os resultados brutos estão organizados em um repositório do projeto, o que permite reproduzir qualquer execução individual a partir de sua semente.

\begin{figure}[htb]
\centering
\begin{tikzpicture}[
  font=\small,
  etapa/.style={draw, text width=2.5cm, minimum height=1.0cm, align=center},
  sub/.style={draw, text width=2.5cm, align=center, font=\scriptsize, inner sep=3pt},
  seta/.style={-{Stealth[length=2.5mm]}, thick}
]
\node[etapa] (base) {Base de dados};
\node[etapa, right=0.7cm of base] (treino) {Treinamento\\com defesa};
\node[etapa, right=0.7cm of treino] (modelo) {Modelo\\treinado};
\node[etapa, right=0.7cm of modelo] (ataque) {Ataque iDLG};

\node[sub, below=0.9cm of base] (bases) {MNIST\\Fashion-MNIST\\CIFAR-10\\4.000 treino\\1.000 teste};
\node[sub, below=0.9cm of treino] (defesas) {MLP, CNN, ViT\\lotes de 1 e 32\\sem defesa\\Top-$k$, Random-$k$\\recorte e ruído\\QSGD};
\node[sub, below=0.9cm of modelo] (util) {acurácia de teste\\(utilidade)};
\node[sub, below=0.9cm of ataque] (metricas) {MSE de reconstrução\\acerto de rótulo\\(privacidade)};

\draw[seta] (base) -- (treino);
\draw[seta] (treino) -- (modelo);
\draw[seta] (modelo) -- (ataque);
\draw[seta] (base) -- (bases);
\draw[seta] (treino) -- (defesas);
\draw[seta] (modelo) -- (util);
\draw[seta] (ataque) -- (metricas);
\end{tikzpicture}
\caption{Fluxograma da metodologia.}
\label{fig:fluxograma}
\end{figure}

\section{Análise Estatística e Critério de Decisão} \label{sec:analise}

As comparações foram organizadas em dois grupos. São confirmatórias as comparações entre cada defesa e o cenário sem defesa, dentro de uma mesma base e arquitetura, que testam diretamente as hipóteses H1 e H2 do Capítulo~1. São exploratórias as comparações entre arquiteturas, entre bases e entre tamanhos de lote, que orientam a discussão da hipótese H3 sem suporte confirmatório.

Para a utilidade, a unidade de análise é a execução de treinamento: cada combinação dispõe de 30 observações independentes, uma por semente. Para a privacidade, a unidade de análise é a imagem atacada: cada combinação dispõe de até 30 observações, obtidas sobre um único modelo treinado. Essa assimetria tem uma consequência importante. Como as 30 reconstruções de uma mesma configuração compartilham o mesmo modelo e o mesmo sorteio da defesa, elas não são inteiramente independentes, e a variabilidade entre treinamentos não é capturada. Os testes sobre o erro de reconstrução devem, portanto, ser lidos como evidência sobre aquele modelo específico, e não como estimativa da variabilidade de toda a população de modelos.

As diferenças foram avaliadas pelo teste $t$ de Welch, conforme a Equação~\ref{eq:welch}, com nível de significância $\alpha = 0{,}05$, e sua magnitude foi quantificada pelo $d$ de Cohen, conforme a Equação~\ref{eq:cohen}. Como são realizadas 36 comparações confirmatórias, quatro defesas em nove combinações de base e arquitetura, a probabilidade de ao menos um falso positivo seria elevada sem ajuste. Por isso, reportam-se o valor-$p$ bruto e o valor corrigido pelo método de Holm, e a significância é declarada apenas quando o valor corrigido permanece abaixo de $\alpha$. O tamanho de efeito é reportado em todos os casos, pois diferenças significativas podem ser pequenas e diferenças não significativas podem ser grandes com amostras reduzidas.

O número de 30 observações por configuração foi escolhido por ser suficiente para detectar, com o teste de Welch e $\alpha = 0{,}05$, efeitos grandes na escala do $d$ de Cohen, da ordem de $0{,}8$ ou mais, com poder próximo de $0{,}8$. Efeitos pequenos ou moderados podem, portanto, passar despercebidos, o que reforça a necessidade de interpretar resultados não significativos como ausência de evidência, e não como evidência de ausência.

Por fim, o critério para identificar as configurações com melhor equilíbrio entre privacidade e utilidade foi definido antes da análise dos resultados. Considera-se preferível toda configuração não dominada, isto é, aquela para a qual nenhuma outra apresente simultaneamente erro de reconstrução maior e acurácia maior, dentro da mesma base e arquitetura. Entre as configurações não dominadas, destacam-se as que preservam ao menos 90\% da acurácia obtida sem defesa e, ao mesmo tempo, aumentam o erro de reconstrução de forma estatisticamente sustentada após a correção de Holm.

\section{Ameaças à Validade e Decisões de Implementação} \label{sec:ameacas-validade}

Esta seção reúne as decisões que limitam o alcance das conclusões, separadas entre as que afetam a validade interna, isto é, a correção das comparações feitas, e as que afetam a validade externa, isto é, a generalização para outros cenários.

Quanto à validade interna, três pontos merecem destaque. O primeiro é que o mecanismo de recorte e ruído foi aplicado ao gradiente médio do lote, e não ao gradiente de cada exemplo, e o ruído não foi dividido pelo tamanho do lote; as duas formulações coincidem apenas quando o lote tem uma amostra. Como também não se utilizou nenhum contador de orçamento de privacidade, esse braço do experimento não implementa privacidade diferencial e não autoriza nenhuma afirmação sobre garantias $(\varepsilon, \delta)$. Soma-se a isso que, com $\sigma C = 1$ em cada uma das centenas de milhares de coordenadas, o ruído supera em várias ordens de grandeza o gradiente recortado, de modo que os resultados desse braço descrevem um regime de ruído dominante, e não o custo da privacidade diferencial em orçamentos usuais.

O segundo ponto é o cálculo do erro de reconstrução. No CIFAR-10, os valores desnormalizados são limitados ao intervalo $[0, 1]$ antes do cálculo, o que não ocorre no MNIST e no Fashion-MNIST. Nessas duas bases, o erro pode ultrapassar $1$, e os valores não devem ser comparados diretamente com os do CIFAR-10; as comparações entre defesas são sempre feitas dentro de uma mesma base e arquitetura.

O terceiro é que o gradiente atacado não corresponde exatamente à mensagem que o cliente enviaria. O treinamento utiliza lotes de uma ou de 32 amostras, enquanto o ataque observa o gradiente de uma única imagem; além disso, as métricas de privacidade foram obtidas sobre um único modelo por configuração, sem replicação entre sementes. A combinação dessas duas escolhas significa que os resultados de privacidade representam o vazamento máximo observável em uma amostra isolada, em um modelo específico, e não a média esperada sobre as rodadas de um treinamento federado.

Quanto à validade externa, o protocolo simula um único cliente, sem agregação, enquanto o aprendizado federado real envolve múltiplos participantes e muitas rodadas. A agregação tende a reduzir o vazamento, de modo que as medidas aqui obtidas são conservadoras do ponto de vista da privacidade; já as medidas de utilidade não se transferem diretamente, pois a acurácia sob FedAvg depende do número de rodadas e da distribuição dos dados entre clientes. No mesmo sentido, o treinamento por uma única época sobre 4.000 imagens mede o custo da defesa no início do aprendizado, e não seu custo assintótico, o que pode favorecer ou penalizar de forma desigual defesas que apenas retardam a convergência.

Por fim, os parâmetros das defesas foram mantidos fixos, com taxa de retenção de 10\%, desvio de ruído unitário e 16 níveis de quantização. O efeito da família de defesa está, portanto, confundido com o efeito da configuração escolhida, e as conclusões valem para esses ajustes específicos, não para as famílias em geral. A varredura desses parâmetros, que produziria curvas completas de compromisso entre privacidade e utilidade, e a adoção de um adversário adaptativo, capaz de incorporar a defesa ao ataque, permanecem como extensões naturais do estudo. Registra-se ainda que, em 21 das 1.350 reconstruções, a otimização divergiu e produziu valores indefinidos; essas tentativas foram excluídas do cálculo do erro de reconstrução, mas mantidas na acurácia de inferência de rótulo, que é obtida antes da otimização.
